\chapter{Pendahuluan}\label{ch:introduction}
    
    \section{Latar Belakang Masalah}\label{sec:background}

        % (KASIH COMMENT SEBELUM `BEGIN` DAN SETELAH `END`)
        % Ilustrasi singkat menggambarkan metode serupa dengan penelitian Tugas Akhir
        \begin{figure}
            \centering
            \includegraphics[width=0.5\linewidth]{example-image-a}
            \caption{Caption}
            \label{fig:ilustrasi model}
        \end{figure}
        %

    \section{Rumusan Masalah}\label{sec:problem}

        % Rumusan Masalah pada Penelitian Tugas Akhir
        \begin{enumerate}
            \item 
        \end{enumerate}
        % 

    \section{Batasan Masalah}\label{sec:limitation}

        % Batasan Masalah pada Penelitian Tugas Akhir
        \begin{enumerate}
            \item 
        \end{enumerate}
        % 

    \section{Tujuan}\label{sec:purpose}

        % Tujuan Penelitian Tugas Akhir
        \begin{enumerate}
            \item 
        \end{enumerate}
        % 

    \section{Manfaat}\label{sec:benefit}

        % Manfaat Penelitian Tugas Akhir
        \begin{enumerate}
            \item 
        \end{enumerate}
        % 